%=====================================================================
\section{Geometry and Gauge from the Gram Matrix}
\label{ch:geometry}
%=====================================================================

Chapters~\ref{ch:whyC}--\ref{ch:whyd5} established $\psi_i \in \mathbb{C}^5$ with $\|\psi_i\| = 1$. We now extract the full content of general relativity and the Standard Model gauge structure from the Gram matrix $G_{ij} = \langle\psi_i|\psi_j\rangle$.

\subsection{The Gram matrix and its properties}

The inner products between all $N$ points define the Gram matrix:
\begin{equation}
G_{ij} = \langle\psi_i|\psi_j\rangle = \sum_{a=0}^{4} \psi^*_{i,a}\,\psi_{j,a}.
\end{equation}

This $N \times N$ matrix has the following properties, all of which are \emph{theorems} (consequences of the inner product structure), not assumptions:

\begin{enumerate}[label=(\alph*)]
\item \textbf{Hermitian:} $G^\dagger = G$, since $G_{ji} = G^*_{ij}$.
\item \textbf{Positive semidefinite:} $\mathbf{c}^\dagger G\,\mathbf{c} = \|\Psi^\dagger\mathbf{c}\|^2 \geq 0$ for all $\mathbf{c} \in \mathbb{C}^N$.
\item \textbf{Rank $\leq 5$:} $G = \Psi\Psi^\dagger$ where $\Psi$ is $N \times 5$, so $\text{rank}(G) \leq 5$.
\item \textbf{Unit diagonal:} $G_{ii} = \|\psi_i\|^2 = 1$ (point democracy).
\item \textbf{Bounded:} $|G_{ij}| \leq 1$ (Cauchy--Schwarz).
\end{enumerate}

Each element $G_{ij}$ is a complex number carrying $2d = 10$ real parameters (a sum of $d = 5$ complex products $\psi^*_{i,a}\psi_{j,a}$). The number of independent components of a 4-dimensional symmetric metric tensor $g_{\mu\nu}$ is $d_\text{st}(d_\text{st}+1)/2 = 10$. This coincidence is exact: one edge of $G$ carries precisely the information content of one 4D metric.

\subsection{Polar decomposition: geometry and gauge}

Each off-diagonal element decomposes into magnitude and phase:
\begin{equation}
G_{ij} = |G_{ij}|\,e^{i\phi_{ij}}, \qquad |G_{ij}| \in [0, 1], \quad \phi_{ij} \in (-\pi, \pi].
\end{equation}

Define the geometric and gauge variables:
\begin{align}
W_{ij} &= \frac{|G_{ij}|^2}{d} \qquad\text{(real, symmetric: } W_{ij} = W_{ji}\text{)}, \label{eq:W}\\
\phi_{ij} &= \arg(G_{ij}) \qquad\text{(real, antisymmetric: } \phi_{ji} = -\phi_{ij}\text{)}. \label{eq:phi}
\end{align}

\textbf{$W$ encodes geometry} (Sec.~\ref{sec:metric}). \textbf{$\phi$ encodes gauge fields} (Sec.~\ref{sec:gauge_connection}).

This separation is the lattice version of the statement: ``gravity is the symmetric part of the connection; gauge forces are the antisymmetric part.''

\subsection{Metric and distance}
\label{sec:metric}

The geodesic distance on $\mathbb{CP}^4$ between points $i$ and $j$ is the Fubini--Study distance:
\begin{equation}
d^2_\text{FS}(i, j) = \arccos^2|G_{ij}|.
\end{equation}

For nearby points ($|G_{ij}| \approx 1$), this linearizes to a metric:
\begin{equation}
ds^2_{ij} = 1 - d\,W_{ij}.
\end{equation}

The interval is spacelike ($ds^2 > 0$) when $W_{ij} < 1/d$, and null ($ds^2 = 0$) when $|G_{ij}| = 1$ (identical points). No two distinct points can have $ds^2 \leq 0$ for generic $\psi$, which eliminates closed timelike curves.

\subsection{Simplicial structure}
\label{sec:simplex}

The rank constraint $\text{rank}(G) \leq 5$ means at most 5 points are linearly independent. Any set of 5 generic points in $\mathbb{CP}^4$ forms a \emph{4-simplex} (pentachoron). The geometry is naturally simplicial: Regge calculus emerges without being imposed.

A 4-simplex has:
\begin{equation}
\binom{5}{1} = 5 \text{ vertices}, \quad
\binom{5}{2} = 10 \text{ edges}, \quad
\binom{5}{3} = 10 \text{ hinges (triangles)}, \quad
\binom{5}{4} = 5 \text{ tetrahedra}, \quad
\binom{5}{5} = 1 \text{ 4-simplex}.
\end{equation}

\subsection{The simplex as data bus: 10 edges = 10 metric components}
\label{sec:data_bus}

The 4-simplex has $\binom{5}{2} = 10$ edges. A 4-dimensional symmetric metric tensor $g_{\mu\nu}$ has $d_\text{st}(d_\text{st}+1)/2 = 10$ independent components. This is not a coincidence: each edge of the simplex maps to one component of the metric tensor via the Regge correspondence.

Each edge carries \emph{two} data channels from the full inner product $G_{ij} = |G_{ij}|\,e^{i\phi_{ij}}$:
\begin{itemize}
\item $|G_{ij}|^2 \to W_{ij}$: the metric component (\textbf{gravity}),
\item $\phi_{ij} = \arg(G_{ij})$: the gauge connection (\textbf{forces}).
\end{itemize}

Under the causal $(n_S, n_T) = (3, 2)$ split, the 10 edges decompose as:

\begin{center}
\begin{tabular}{lccl}
\hline
Edge type & Count & Formula & Metric role \\
\hline
Temporal--Temporal & $\binom{2}{2} = 1$ & $W$ within time & $g_{00}$ (lapse) \\
Spatial--Spatial & $\binom{3}{2} = 3$ & $W$ within space & $g_{ij},\;i\neq j$ (spatial off-diagonal) \\
Spatial--Temporal & $2\times 3 = 6$ & $W$ across sectors & $g_{0i}$ (shift) + diagonal \\
\hline
\textbf{Total} & \textbf{10} & & $g_{\mu\nu}$ fully determined \\
\hline
\end{tabular}
\end{center}

The simplex is a minimal processor: 5 complex-number registers (vertices) connected by 10 pairwise comparisons (edges) whose magnitudes encode local geometry and whose phases encode gauge fields. The simplex interior carries no information --- it is the irreducible ``blind spot'' that defines the minimum resolution unit.

\subsection{Lorentz signature from unitarity}
\label{sec:Lorentz}

The metric $ds^2_{ij} = 1 - d\,W_{ij}$ is positive-definite (Riemannian). All distances are spacelike. Where does the Lorentzian signature $(-,+,+,+)$ come from?

\begin{theorem}[Lorentz signature is derived]
Unitarity of time evolution forces Lorentzian signature.
\end{theorem}

\begin{proof}
The argument has four steps.

\medskip
\noindent\textbf{Step 1: Unitarity selects a time direction.}
Time evolution between spatial slices $\Sigma_t$ and $\Sigma_{t+1}$ is unitary: $U^\dagger U = I$ (a consequence of inner-product preservation). This foliation breaks the directional symmetry of $G$: directions \emph{within} a slice are spatial; directions \emph{between} slices are temporal.

\medskip
\noindent\textbf{Step 2: The factor of $i$ is non-negotiable.}
The unitary propagator $U(\Delta t) = e^{-iH\Delta t}$ requires the imaginary unit. The alternatives fail: $e^{-H\Delta t}$ decays (not unitary), $e^{+H\Delta t}$ grows (not unitary). Only $e^{-iH\Delta t}$ oscillates, conserving probabilities.

\medskip
\noindent\textbf{Step 3: Wick rotation is forced.}
The path integral weight $e^{iS/\hbar}$ relates to the Euclidean weight $e^{-S_E/\hbar}$ via $\tau = it$. This Wick rotation transforms the Euclidean metric $ds^2_E = d\tau^2 + d\mathbf{x}^2$ into the Lorentzian metric $ds^2_L = -dt^2 + d\mathbf{x}^2$: the temporal component flips sign.

\medskip
\noindent\textbf{Step 4: Lattice metric.}
On the lattice, spacelike edges (within a slice) retain $ds^2_\text{space} = 1 - d\,W_{ij} > 0$, while timelike edges (between slices) acquire $ds^2_\text{time} = -(1 - d\,W_{ij}) < 0$. The combined signature is $(-,+,+,+)$.
\end{proof}

\begin{remark}
The light cone corresponds to $ds^2 = 0$, i.e., $|G_{ij}| = 1$: the boundary where two points become indistinguishable. Lorentz invariance $\text{SO}(1,3)$ is the symmetry group preserving this derived signature --- it is a theorem, not a postulate. Special and general relativity follow as the flat and curved limits, respectively.
\end{remark}

\subsection{Hinge area and curvature}

A hinge (triangle) with vertices $\{i, j, k\}$ has a $3\times 3$ Gram sub-matrix:
\begin{equation}
G_h = \begin{pmatrix} 1 & G_{ij} & G_{ik} \\ G_{ji} & 1 & G_{jk} \\ G_{ki} & G_{kj} & 1 \end{pmatrix}.
\end{equation}

Its area in $\mathbb{CP}^4$ is:
\begin{equation}
A_h = \sqrt{\det(G_h)}.
\end{equation}

Expanding the determinant:
\begin{equation}
\det(G_h) = 1 - |G_{ij}|^2 - |G_{jk}|^2 - |G_{ki}|^2 + 2|G_{ij}||G_{jk}||G_{ki}|\cos\Phi_h
\label{eq:detGh}
\end{equation}
where $\Phi_h = \phi_{ij} + \phi_{jk} + \phi_{ki}$ is the \emph{holonomy} around the triangle --- the gauge-invariant lattice field strength (Sec.~\ref{sec:holonomy}).

The dihedral angle at hinge $h$ is computed from the two tetrahedra sharing $h$, and the deficit angle $\delta_h = 2\pi - \sum_{\sigma \ni h} \theta_\sigma$ measures curvature. The Regge action is:
\begin{equation}
\boxed{S = \sum_h A_h\,\delta_h}
\label{eq:Regge}
\end{equation}
which converges to the Einstein--Hilbert action $\int R\sqrt{g}\,d^4x$ in the continuum limit \cite{Regge}.

\begin{remark}[No coupling constant in the action]
The action (\ref{eq:Regge}) contains no adjustable parameters: no $G_\text{Newton}$, no $\hbar$, no $\alpha$. It is pure geometry. All coupling constants emerge from the \emph{structure} of $G$, not from coefficients in the action (see Chapter~\ref{ch:couplings}).
\end{remark}

\subsection{Gauge connection from phase}
\label{sec:gauge_connection}

The phase $\phi_{ij} = \arg\langle\psi_i|\psi_j\rangle$ transforms under local rephasing $\psi_i \to e^{i\alpha_i}\psi_i$ as:
\begin{equation}
\phi_{ij} \to \phi_{ij} + \alpha_j - \alpha_i
\end{equation}
which is precisely the lattice gauge transformation $A_{ij} \to A_{ij} + \partial_j\alpha - \partial_i\alpha$. Thus $\phi_{ij}$ is a \emph{lattice gauge connection}.

\subsection{Holonomy as field strength}
\label{sec:holonomy}

The holonomy around a triangle $\{i, j, k\}$:
\begin{equation}
\Phi_{ijk} = \phi_{ij} + \phi_{jk} + \phi_{ki} = \arg(G_{ij}\,G_{jk}\,G_{ki})
\end{equation}
is gauge-invariant ($\alpha_i$ cancels cyclically). This is the discrete analogue of the field strength $F_{\mu\nu} = \partial_\mu A_\nu - \partial_\nu A_\mu + [A_\mu, A_\nu]$.

From Eq.~(\ref{eq:detGh}), the hinge area $A_h = \sqrt{\det(G_h)}$ depends on $\Phi_h$: geometry and gauge are coupled through the determinant. This coupling is Einstein's equation in discrete form: matter (gauge fields) curves spacetime (hinge areas).

\subsection{The causal split: $(n_S, n_T) = (3, 2)$}

The 5 vertices of the simplex split into spatial and temporal:
\begin{equation}
\underbrace{a,\,b,\,c}_{n_S = 3\;\text{(spatial)}} \quad \underbrace{d,\,e}_{n_T = 2\;\text{(temporal)}}
\end{equation}

This split is unique: chirality of the weak force requires $n_T = 2$. Among all $\text{SU}(n)$, only $\text{SU}(2)$ has a pseudo-real fundamental representation ($\mathbf{2} \cong \bar{\mathbf{2}}$ via the antisymmetric tensor $\epsilon_{ab}$), which is the algebraic condition for chiral fermions. This forces the temporal sector to be $\mathbb{C}^2$, and therefore $n_S = d - n_T = 3$.

\subsection{Speed of light}

The lattice speed of light is the ratio of temporal to spatial vertex densities per dimension:
\begin{equation}
\boxed{c = \frac{n_T / d_T}{n_S / d_S} = \frac{2/1}{3/3} = 2}
\label{eq:speed_of_light}
\end{equation}
where $d_T = 1$ (time has 1 macroscopic dimension) and $d_S = 3$ (space has 3).

Every factor of $\sqrt{2}$ in the Standard Model --- $m = yv/\sqrt{2}$, the Higgs normalization, etc.\ --- is $\sqrt{c}$. The speed of light is not a fundamental constant but a \emph{lattice ratio}: time has twice the vertex density of space, per dimension.

\begin{remark}
The constancy of $c$ follows from its being a ratio of integers. It cannot vary, because $n_S = 3$ and $n_T = 2$ are fixed by the uniqueness of $d = 5$ (Theorem~\ref{thm:d5}).
\end{remark}

\subsection{Gauge group from chirality}

The full symmetry of $\mathbb{C}^5$ is $\text{SU}(5)$. The causal split breaks this:
\begin{equation}
\text{SU}(5) \;\to\; \text{SU}(3) \times \text{SU}(2) \times \text{U}(1)
\end{equation}

This is the Georgi--Glashow decomposition \cite{GeorgiGlashow}, but here it follows from the axiom rather than being postulated. The three factors are:
\begin{itemize}
\item $\text{SU}(3)$: internal rotations of $\mathbb{C}^3$ (strong force, color),
\item $\text{SU}(2)$: internal rotations of $\mathbb{C}^2$ (weak force, isospin),
\item $\text{U}(1)$: relative phase between $\mathbb{C}^3$ and $\mathbb{C}^2$ (electromagnetic force).
\end{itemize}

\subsection{Weinberg angle}

In $\text{SU}(5)$, the weak mixing angle at the unification scale is fixed by the trace condition on the generators:
\begin{equation}
\sin^2\theta_W = \frac{\text{Tr}(T_3^2)}{\text{Tr}(Q^2)} = \frac{3}{8}
\end{equation}
where $T_3$ and $Q$ are evaluated in the fundamental $\mathbf{5}$ representation. This is the standard Georgi--Glashow result, but here it is a consequence of $\mathbb{C}^5 = \mathbb{C}^3 \oplus \mathbb{C}^2$.

\subsection{Hinge classification and force hierarchy}

The 10 hinges of the 4-simplex classify by their vertex content:

\begin{center}
\begin{tabular}{lccl}
\hline
Type & Count & $\det(G_h)$ & Force \\
\hline
SSS (3 spatial) & $\binom{3}{3}\binom{2}{0} = 1$ & $> 0$ & Strong \\
SST (2 spatial, 1 temporal) & $\binom{3}{2}\binom{2}{1} = 6$ & $> 0$ & EM / Weak \\
STT (1 spatial, 2 temporal) & $\binom{3}{1}\binom{2}{2} = 3$ & $= 0$ always & Classical \\
TTT (3 temporal) & $\binom{2}{3} = 0$ & --- & Impossible \\
\hline
\end{tabular}
\end{center}

The STT hinges have $\det(G_h) = 0$ because $n_T = 2 < 3$: two temporal vertices cannot span a nondegenerate triangle. TTT is impossible since $\binom{2}{3} = 0$.

This gives a natural force hierarchy: 7 quantum hinges (SSS + SST) mediate the strong and electroweak forces, while STT hinges are classical ($\hbar = 0$). The weak force has no dedicated hinge (TTT $= 0$) and must ``borrow'' from the SST sector, providing a geometric explanation for its weakness.

\subsection{Singularity resolution}

\begin{proposition}
No two distinct points can have $ds^2 = 0$.
\end{proposition}

\begin{proof}
$ds^2 = 0$ requires $|G_{ij}| = 1$, i.e., $\psi_i = e^{i\theta}\psi_j$, which means $i$ and $j$ represent the same physical point in $\mathbb{CP}^4$. Distinct points always have $ds^2 > 0$.
\end{proof}

This eliminates singularities: gravitational collapse drives $\psi_i \to \psi_j$ (alignment), giving $\det(G_h) \to 0$, hence $\hbar_\text{eff} \to 0$ (Chapter~\ref{ch:hbar}). The region becomes classical vacuum, not a singularity. The ``center'' of a black hole is where $\psi$ values align --- it is vacuum, not infinite density.
